\section{Security context agent}
\label{sec:agent}

\paragraph{Problem and inputs.}
Earlier rounds supplied acquisition with researcher-written security
perspectives; a practitioner who can write them can also write the guards.
The agent removes that input and assembles $\mathrm{SC}(R,\Delta)$ from a
read-only snapshot of the target repository (for reuse tasks also the donor),
the one-line feature description given to the generator, and the ontology of
Section~\ref{sec:terminology} as an output schema with its catalogue
vocabulary (CWE, CAPEC, ASVS). It receives no check, fixture, threshold,
earlier security context, generated code, audit or test outcome, executes no
repository code, and runs under a fixed, recorded tool budget, model and
sampling settings.

\paragraph{Code model.}
The harness parses every source file into a code model (types, methods,
constructors, fields, call edges, sink tags). Every statement is anchored to
a code-model symbol, a file range or both.

\paragraph{Angles.}
The agent runs an inspection loop with shell tools~\citep{yang2024} and
produces one typed document. Three angles:
\begin{enumerate}\setlength\itemsep{0pt}
\item \emph{Data flow}: trace untrusted data from the touched code across
  trust boundaries to sinks, enumerate threats per boundary (STRIDE), map
  them to weakness types with a reachability judgement, and place requirement
  and control where the data first arrives.
\item \emph{Requirements}: name assets and their properties, then constrain
  every feature operation with a requirement, enforcement point and control
  with failure behaviour.
\item \emph{Catalogue}: decide with evidence, per entry of a published
  weakness catalogue~\citep{cwetop25}, whether it applies to this change;
  derive requirements and controls only for applicable entries.
\end{enumerate}

\paragraph{Output.}
Each statement carries a kind, a basis (observed, task, reasoned), catalogue
identifiers and references to inspected excerpts (Table~\ref{tab:terms}).
The harness, not the model, resolves references to file and line ranges and
hashes the cited text; an unknown symbol with a verifiable file range is
kept as the range and the correction counted; an unresolvable reference
invalidates the statement. The validated document is stored as a typed graph
(Table~\ref{tab:terms}) and rendered verbatim, in a fixed order, as an insert
after the unchanged task; controls receive no insert.

\begin{table}[t]
\centering
\footnotesize
\setlength{\tabcolsep}{4pt}
\caption{Terminology, schema and graph. Records before schema v8 label observations \emph{security\_property}.}
\label{tab:terms}
\resizebox{\textwidth}{!}{\begin{tabular}{@{}lll@{}}
\toprule
Term (Section~\ref{sec:terminology}) & Schema element & Graph \\
\midrule
Asset, security property & \texttt{assets[].name}, \texttt{.property} & asset node, \emph{located\_at} symbol \\
Trust boundary, entry point & \texttt{boundaries[]}, \texttt{.entry\_point} & boundary node, \emph{crosses} symbol \\
Attack surface & boundaries with \texttt{entry\_point} & entry-point boundary nodes \\
Threat & \texttt{items[].threat} (STRIDE) & node attribute \\
Attack pattern, weakness type & \texttt{items[].capec}, \texttt{.cwe} & node attributes \\
Observed fact & kind \texttt{observation}, basis observed & node, \emph{anchored\_at} symbol \\
Existing or change risk & kinds \texttt{existing\_risk}, \texttt{change\_risk} & node, \emph{weakens} observation \\
Security requirement & kind \texttt{requirement} & node, \emph{mitigates} risk \\
Control, enforcement point, failure behaviour & kind \texttt{control}, \texttt{.enforcement\_point}, \texttt{.failure\_behavior} & node, \emph{satisfies} requirement, \emph{enforced\_at} symbol \\
Verification & kind \texttt{verification}, \texttt{.verification} & node, \emph{verifies} \\
Unknown & kind \texttt{unknown} & node without anchors \\
Code element & anchor symbol, file, line range & symbol and file nodes, \emph{calls} edges \\
\bottomrule
\end{tabular}}
\end{table}

\paragraph{Delivery modes.}
\emph{Single-shot} delivery is the protocol of Section~\ref{sec:design}; in
\emph{agentic} delivery the generator may search and read the repository
before submitting. Sidecars: none; static (the full insert); adaptive (no upfront insert; after each read and before each
submission the harness injects the graph \emph{slice} for the code touched
so far: the statements anchored at the symbols read or edited or at their
call neighbours, with their requirements and controls, once each);
gatekeeper (injects nothing; when a submission edits an enforcement point, a
one-turn judge, the same model, recorded, rejects it with the violated
statements only if the edit clearly violates an anchored control). The judge
sees the edit and the statements, never test outcomes.

\paragraph{Provenance.}
A record freezes instruction, ontology and snapshot fingerprints, every tool
call and model turn, the validated output and its insert, with hashes, before
code generation. Failed or degenerate acquisitions are kept and
used, none regenerated: in I09 one catalogue acquisition inspected nothing
(wrong path syntax in its tool calls); its insert was used unchanged and
reported as uninspected.

\paragraph{Ablations.}
Two arms bracket the angles: \emph{task-only}, a single sentence without
ontology, code model or procedure, and a flat \emph{catalogue} arm adding the
CWE Top~25 list~\citep{cwetop25} without procedure or code model. Differences
between arms confound instruction, generated content and insert length; the
design does not isolate wording.

\paragraph{Relation to prior work.}
Prior work evaluates code-model security in isolated
scenarios~\citep{pearce2022,siddiq2022,hajipour2024}, hardens the
model~\citep{he2023} or varies generic prompt wording~\citep{tony2025}; user
studies find assistance without security cues degrades
security~\citep{perry2023}. Here repository-specific security context is
derived from evidence before generation, in a corpus where flaws propagate by
clone-and-own reuse~\citep{krueger2018,dubinsky2013}.
